% Generated draft: requires mathematical and source review.
\documentclass[11pt]{article}
\usepackage[margin=25mm]{geometry}
\usepackage{fontspec}
\setmainfont{Latin Modern Roman}
\usepackage{amsmath,amssymb,mathtools,unicode-math}
\setmathfont{Latin Modern Math}
\usepackage{enumitem,needspace,xurl,hyperref}
\setlength{\parindent}{0pt}
\setlength{\parskip}{5pt}
\setlength{\emergencystretch}{3em}
\newcommand{\N}{\mathbb{N}}
\newcommand{\Z}{\mathbb{Z}}
\newcommand{\Q}{\mathbb{Q}}
\newcommand{\R}{\mathbb{R}}
\newcommand{\C}{\mathbb{C}}
\newcommand{\F}{\mathbb{F}}
\newcommand{\A}{\mathbb{A}}
\newcommand{\PP}{\mathbb{P}}
\newcommand{\E}{\mathbb{E}}
\newcommand{\eps}{\varepsilon}
\newcommand{\dd}{\,\mathrm{d}}
\begin{document}
\section*{UnsolvedMath ID: 30003471}
\textbf{Source problem code:} catalogue entry
\subsection{Definitions and mathematical statement}
\textit{Provisional conventions and unresolved hypotheses.} A convex polytope is the convex hull of a finite set of points. To define the usual dihedral angles in this draft, suppose that each polytope has affine dimension three. Let \(E(P)\) denote the set of edges of \(P\). A combinatorial equivalence is an inclusion-preserving isomorphism \(\Phi\) between the face lattices of two polytopes; it induces a bijection between their edge sets. For an edge \(e\in E(P)\), let \(\alpha_P(e)\in(0,\pi)\) be the interior dihedral angle between the two facets incident with \(e\), measured inside \(P\).

\textit{Literal comparison question.} Given combinatorially equivalent convex polytopes \(P_1,P_2\subset\R^3\) and a correspondence \(\Phi\), must one have
\[
\left(\exists e\in E(P_1):\ \alpha_{P_1}(e)\leq\alpha_{P_2}(\Phi(e))\right)
\quad\lor\quad
\left(\forall e\in E(P_1):\ \alpha_{P_1}(e)=\alpha_{P_2}(\Phi(e))\right)?
\]
The source asks this for arbitrary such pairs of polytopes, but does not specify whether the correspondence is prescribed, universally quantified, or existentially chosen. It also does not explicitly require affine dimension three or specify interior versus exterior dihedral angles. The conventions above therefore require confirmation.

The displayed disjunction preserves the supplied non-strict comparison and equality alternative. Under the stated conventions, the equality alternative already implies the first alternative. The supplied record does not establish whether a strict comparison outside the all-equal case was intended; no such strengthening is asserted here.
\subsection{Short English statement}
The question compares dihedral angles at corresponding edges of two combinatorially equivalent convex polytopes in \(\R^3\). It asks whether some angle of the first is no greater than its partner in the second, with an all-equal alternative. Full dimensionality, the angle convention, and the quantifier on the correspondence require confirmation; the equality alternative is redundant with the supplied non-strict inequality.
\subsection{Sources}
\begin{itemize}
\item Local UnsolvedMath registry, permanent ID \(30003471\), source problem code \texttt{catalogue entry}, title: Dihedral Angle Comparison for Combinatorially Equivalent Polytopes. The supplied record belongs to the Oberwolfach Reports Open Problems collection and lists proposed year 2017; no proposer is named.
\item Supplied report citation: \emph{Discrete Geometry (2018)}, Oberwolfach Reports, DOI 10.4171/owr/2017/19, pp.\ 1135--1205. \url{https://doi.org/10.4171/owr/2017/19}. The background describes these reports as records of workshops and their open-problem discussions. The citation's 2018 designation and the record's 2017 proposed year are retained without reconciliation.
\item Supplied literature reference: Akopyan and Karasev on polytope angles, \url{https://arxiv.org/abs/1505.05263}. The supplied assessment identifies the assertion as their Conjecture 5.1; this attribution and identification have not been independently verified.
\item Supplied literature assessment dated exactly \texttt{2026-08-21}: the record labels the problem open, reports known simplex cases and related angle inequalities, and states that no general proof or counterexample for combinatorially equivalent three-dimensional convex polytopes was found. The supplied checked-source links are \url{https://arxiv.org/abs/1505.05263} and \url{https://doi.org/10.4171/owr/2017/19}. This is supplied dated curation derived from the OWR record, not an independently verified literature assessment.
\end{itemize}
\end{document}
